\chapter{Mechanization Status and Open Obligations}
\index{mechanized proof!status|textbf}

\label{ch:mechanization-status}

\chapterstatus{proof contracts and proof strategies are presented with explicit status; no proposed physics proof package is represented as complete or reviewed.}


\section{Mechanization status}
\index{Lean}
\index{Isabelle}
\index{Rocq}


The current prover-neutral \texttt{PRF-05} target set contains nine frozen
contracts:

\begin{center}
\small
\begin{tabular}{@{}lll@{}}
\toprule
Contract & Subject & Current backend evidence \\
\midrule
\texttt{05.01} & Projector invariance & Lean checked \\
\texttt{05.02} & Compression covariance & Unencoded \\
\texttt{05.03} & Pullback covariance & Unencoded \\
\texttt{05.04} & Aligned subtraction & Unencoded \\
\texttt{05.05a} & Frobenius-norm invariance & Unencoded \\
\texttt{05.05b} & Operator-norm invariance & Unencoded \\
\texttt{05.06} & Equivariant path residual & Unencoded \\
\texttt{05.07} & Gauge--truncation counterexample & Unencoded \\
\texttt{05.08} & Shell-Frobenius invariance & Unencoded \\
\bottomrule
\end{tabular}
\end{center}

Every Isabelle and Rocq target is unencoded. Eight Lean targets are unencoded.
No contract is cross-checked. The package remains \texttt{proposed}, while
\texttt{PRF-05.01} is individually \texttt{in development} because the
multi-backend completion and review criteria have not been met.

\citationtodo{medium priority}{Check and add scholarly system references for
Lean 4, mathlib, Isabelle/HOL, and Rocq. Provisional keys:
\texttt{demouraUllrich2021}, \texttt{mathlib2020},
\texttt{nipkowPaulsonWenzel2002}, and \texttt{bertotCasteran2004}. The table's
checked and unencoded statuses require durable internal provenance; system
citations do not establish those statuses.}

\section{Novelty and evidence boundary}
\index{proof!evidence boundary}


The finite-dimensional projector, compression, pullback, subtraction, and norm
identities are standard linear algebra. The project does not claim those
identities as new mathematics. A possible contribution is their
application-specific organization into one prover-neutral contract set,
independent Lean--Isabelle--Rocq encodings, semantic conformance review, dual
common-space impurity extraction, and composition with later locality,
continuum, and compatibility arguments. Priority over prior formalization has
not been established.

Potentially substantive mathematical results lie in the assumptions and bounds
that connect the elementary layer to the physical reduction: existence and
regularity of an adequate identification, gauge-constrained locality,
quantitative excluded-space control, well-posed crossover criteria,
operator-to-observable bounds, certified model-class incompatibility, and
atomistic-to-continuum consistency. Every such result remains proposed.

Proof establishes only the stated conclusion from the stated assumptions.
Numerical verification must still test implementations and bound usefulness;
scientific validation must still test the physical model for a declared use;
and human review must still decide whether the assumptions and interpretation
are acceptable.
